%!TEX root = ../main_thesis.tex

\chapter{Dalla libreria alla piattaforma web}
\label{cap:piattaforma}

\todo[inline]{Capitolo di supporto: mostra come il motore (libreria Rust) diventi una libreria riusabile e una piattaforma web didattica. Mantenere il taglio sintetico, coerente con il focus sul motore.}

\section{Compilazione in WebAssembly e la libreria TypeScript}
\label{sec:wasm}

La libreria Rust viene compilata in WebAssembly tramite \texttt{wasm-pack} ed
esposta a JavaScript/TypeScript attraverso \texttt{serde-wasm-bindgen}. Sopra
di essa è costruita la libreria \texttt{@specy/rooc}, che offre un'API
type-safe.

\todo[inline]{Descrivere il processo di build (Rust $\rightarrow$ WASM $\rightarrow$ pacchetto npm), la marshalling dei dati e l'API pubblica di \texttt{ts-lib}. Riferimenti: \texttt{src/lib.rs}, \texttt{src/pipe/pipe\_wasm\_runner.rs}, \texttt{ts-lib/}.}

\section{L'IDE web e l'LSP nel browser}
\label{sec:ide-web}

La piattaforma web è un'applicazione SvelteKit con editor Monaco. Un Language
Server Protocol (LSP) parziale, eseguito interamente nel browser, fornisce
evidenziazione della sintassi, hover dei tipi, completamento ed evidenziazione
degli errori.

\todo[inline]{Descrivere l'IDE: editor Monaco con linguaggio ROOC custom, l'LSP nel browser (diagnostica, hover, completamento, formattazione), il rendering LaTeX con KaTeX e la persistenza locale (IndexedDB). Riferimenti: \texttt{frontend/src/lib/Monaco.ts}, \texttt{frontend/src/lib/Rooc/}.}

\section{Valore didattico}
\label{sec:valore-didattico}

\todo[inline]{Argomentare il valore didattico: visualizzazione passo-passo della risoluzione, esecuzione di funzioni utente in sandbox, documentazione interattiva e condivisione dei modelli. Collegare all'astrazione a pipe (Cap.~\ref{cap:architettura}).}
